% Options for packages loaded elsewhere
\PassOptionsToPackage{unicode}{hyperref}
\PassOptionsToPackage{hyphens}{url}
%
\documentclass[
  12pt,
  a4paper,
]{article}
\usepackage{amsmath,amssymb}
\usepackage{iftex}
\ifPDFTeX
  \usepackage[T1]{fontenc}
  \usepackage[utf8]{inputenc}
  \usepackage{textcomp} % provide euro and other symbols
\else % if luatex or xetex
  \usepackage{unicode-math} % this also loads fontspec
  \defaultfontfeatures{Scale=MatchLowercase}
  \defaultfontfeatures[\rmfamily]{Ligatures=TeX,Scale=1}
\fi
\usepackage{lmodern}
\ifPDFTeX\else
  % xetex/luatex font selection
  \setmainfont[]{DejaVu Serif}
  \setsansfont[]{DejaVu Sans}
\fi
% Use upquote if available, for straight quotes in verbatim environments
\IfFileExists{upquote.sty}{\usepackage{upquote}}{}
\IfFileExists{microtype.sty}{% use microtype if available
  \usepackage[]{microtype}
  \UseMicrotypeSet[protrusion]{basicmath} % disable protrusion for tt fonts
}{}
\makeatletter
\@ifundefined{KOMAClassName}{% if non-KOMA class
  \IfFileExists{parskip.sty}{%
    \usepackage{parskip}
  }{% else
    \setlength{\parindent}{0pt}
    \setlength{\parskip}{6pt plus 2pt minus 1pt}}
}{% if KOMA class
  \KOMAoptions{parskip=half}}
\makeatother
\usepackage{xcolor}
\usepackage[left=15mm,right=15mm,top=20mm,bottom=20mm]{geometry}
\usepackage{longtable,booktabs,array}
\usepackage{calc} % for calculating minipage widths
% Correct order of tables after \paragraph or \subparagraph
\usepackage{etoolbox}
\makeatletter
\patchcmd\longtable{\par}{\if@noskipsec\mbox{}\fi\par}{}{}
\makeatother
% Allow footnotes in longtable head/foot
\IfFileExists{footnotehyper.sty}{\usepackage{footnotehyper}}{\usepackage{footnote}}
\makesavenoteenv{longtable}
\usepackage{graphicx}
\makeatletter
\def\maxwidth{\ifdim\Gin@nat@width>\linewidth\linewidth\else\Gin@nat@width\fi}
\def\maxheight{\ifdim\Gin@nat@height>\textheight\textheight\else\Gin@nat@height\fi}
\makeatother
% Scale images if necessary, so that they will not overflow the page
% margins by default, and it is still possible to overwrite the defaults
% using explicit options in \includegraphics[width, height, ...]{}
\setkeys{Gin}{width=\maxwidth,height=\maxheight,keepaspectratio}
% Set default figure placement to htbp
\makeatletter
\def\fps@figure{htbp}
\makeatother
\setlength{\emergencystretch}{3em} % prevent overfull lines
\providecommand{\tightlist}{%
  \setlength{\itemsep}{0pt}\setlength{\parskip}{0pt}}
\setcounter{secnumdepth}{-\maxdimen} % remove section numbering
\usepackage{sectsty}
\allsectionsfont{\raggedright}
\usepackage{lineno}
\linenumbers
\usepackage{setspace}
\doublespacing
\usepackage{fancyhdr}
\pagestyle{fancy}\fancyhf{}
\fancyhead[L]{\small Lunar and solar tidal responses with optical clocks}
\fancyfoot[C]{\thepage}\setlength{\headheight}{14pt}
\usepackage{float}\floatplacement{figure}{H}
\setlength{\emergencystretch}{2em}
\ifLuaTeX
  \usepackage{selnolig}  % disable illegal ligatures
\fi
\IfFileExists{bookmark.sty}{\usepackage{bookmark}}{\usepackage{hyperref}}
\IfFileExists{xurl.sty}{\usepackage{xurl}}{} % add URL line breaks if available
\urlstyle{same}
\hypersetup{
  hidelinks,
  pdfcreator={LaTeX via pandoc}}

\author{}
\date{}

\begin{document}

\hypertarget{resolving-lunar-and-solar-tidal-responses-with-an-optical-clock-baseline}{%
\section{Resolving lunar and solar tidal responses with an optical clock
baseline}\label{resolving-lunar-and-solar-tidal-responses-with-an-optical-clock-baseline}}

\begingroup\singlespacing

{[}Author list, order and affiliation assignments to be confirmed by the
authors{]}

\textsuperscript{1} Hefei National Laboratory, University of Science and
Technology of China, Hefei, China

\textsuperscript{2} Shanghai Research Center for Quantum Sciences and
CAS Center for Excellence in Quantum Information and Quantum Physics,
University of Science and Technology of China, Shanghai, China

\textsuperscript{3} Hefei National Research Center for Physical Sciences
at the Microscale and School of Physical Sciences, University of Science
and Technology of China, Hefei, China

\textsuperscript{4} Innovation Academy for Precision Measurement Science
and Technology, Chinese Academy of Sciences, Wuhan, China

\textsuperscript{5} Jinan Institute of Quantum Technology and CAS Center
for Excellence in Quantum Information and Quantum Physics, University of
Science and Technology of China, Jinan, China

\textsuperscript{6} Time and Frequency Geodesy Center (TFGC), School of
Geodesy and Geomatics, Wuhan University, Wuhan, China

\endgroup

\textbf{Optical clocks measure tides through changes in their relative
rates, but neighbouring tidal components are difficult to separate when
observations are interrupted and frequency reference levels differ
between segments.\textsuperscript{1,2} Here we use the changing
astronomical phase between lunar and solar tides as an external
reference to resolve neighbouring semidiurnal responses. Clocks in
coastal Shanghai and inland Wuhan accumulated 545.9 hours over 353 days.
Their measured spectrum has local peaks near the principal lunar and
solar frequencies, 0.78 µHz apart. Joint physical-template estimates
give \(1.00\pm0.13\) for the lunar group M₂+N₂ and \(0.93\pm0.36\) for
the neighbouring group S₂+K₂ relative to the combined solid-Earth and
ocean-loading prediction (one statistical standard error). Astronomical
phase thus allows neighbouring tidal responses to be distinguished
across interruptions without requiring a calibrated frequency reference
level shared by all observing segments.}

\clearpage

Earth's gravitational potential varies as the Moon and Sun change
position and as water, air and the solid Earth move. Measuring these
changes provides a way to study mass redistribution and deformation, and
to maintain consistent height and time
references.\textsuperscript{3,4,5} Geodetic surveys and satellite
gravity observations constrain different aspects of this evolving
field.\textsuperscript{4,5} Optical clocks offer a direct measure of
potential differences: gravity changes their relative rates, so a
frequency comparison between distant clocks contains a geophysical
observable.\textsuperscript{1,2} Connecting stable clocks through
optical fibres allows repeated observations of the same endpoints across
hundreds of kilometres.\textsuperscript{6,7}

Laboratory redshift tests and remote optical-clock comparisons have
established measurements of static potential
differences.\textsuperscript{8,9,10} A time-dependent measurement
additionally requires enough precision over the period of the change and
sufficiently stable operation of both clocks and their connection. Tides
are a well-defined target: celestial motion fixes their frequencies and
phases, while deformation and ocean loading determine the signal
expected at each station.\textsuperscript{3} Measuring that signal tests
whether clock comparisons can follow the changing Earth. The same
variation enters the redshift correction needed to compare clocks and
form common time scales.\textsuperscript{11} Coordinated comparisons
already connect optical standards across national
borders.\textsuperscript{12}

The principal lunar and solar semidiurnal tides are a demanding pair of
nearby signals. M₂ repeats every 12.42 h and S₂ every 12 h, giving a
frequency separation of only 0.78 µHz. Over a short comparison their
variations can look much alike. The lunar peaks arrive about 50 minutes
later in solar time each day, so the relative phase of the two tides
changes with the observing date. Measurements on different dates
therefore sample different combinations of the two signals.

An interruption can leave the frequency reference level of a new
observing segment unknown, while its astronomical phase remains
determined by the actual observation time. Each segment still records
the changes within it. Combining these changes at their known
lunar--solar phases supplies separate constraints on the two responses
without first aligning the segments to a common frequency level. This
physical reference motivates retaining the actual UTC times, assigning
each segment an independent constant and estimating the astronomical
templates together.

Here we measure tidal variations of the potential difference between
coastal Shanghai and inland Wuhan and separate neighbouring semidiurnal
responses. Across a 353-day span, 54 valid segments provide 545.9 h of
observations. The two target responses are represented by complete
physical templates for the lunar semidiurnal group M₂+N₂ and the
neighbouring group S₂+K₂. The joint model also retains the approximately
daily tides as an auxiliary physical component. Internal amplitude
ratios and phases remain fixed by the astronomical model. K₂ contains
lunar and solar forcing. The measured spectrum has local peaks near M₂
and S₂, and the joint estimates are consistent with the combined
solid-Earth and ocean-loading prediction.

\hypertarget{a-coastalinland-prediction}{%
\subsubsection{A coastal--inland
prediction}\label{a-coastalinland-prediction}}

The strontium clock in Shanghai and ytterbium clock in Wuhan are 692 km
apart, joined through Hefei by a 1,350-km phase-stabilized fibre link
(Fig. 1a).\textsuperscript{13} The summer observations overlap the
absolute frequency-ratio measurement reported by Hou and colleagues; the
present observable is the time-dependent tidal response. The clock
systems and their performance are described in their published
evaluations.\textsuperscript{14,15} The link extends a transmission
architecture previously demonstrated over 2,067 km of field-deployed
fibre.\textsuperscript{16} We define a positive frequency-ratio
variation as Wuhan Yb running faster relative to Shanghai Sr.~In the
geodetic potential convention used here, the measured fractional
variation is
\(y=-[\delta W_{\mathrm{Wuhan}}-\delta W_{\mathrm{Shanghai}}]/c^2\).

Ocean mass redistribution produces different gravitational responses at
coastal Shanghai and inland Wuhan. Their differential tidal signal
therefore includes an appreciable loading contribution alongside the
solid-Earth tide. The solid response combines the direct astronomical
potential, the Earth's induced potential and displacement of the clocks.
Ocean loading likewise includes water attraction, induced potential and
clock displacement. We combine an IERS-based solid-tide prediction with
FES2022 ocean elevations and a loading kernel containing all three
loading contributions (Methods).\textsuperscript{17,18,19,20}

The predicted M₂ amplitudes are \(3.43\times10^{-18}\) for the solid
contribution and \(1.42\times10^{-18}\) for loading. Their phases differ
by 19.8°, giving a combined amplitude of \(4.78\times10^{-18}\) (Fig.
1b). Loading is about 30\% of the combined M₂ amplitude. On the sampled
record, after removing segment means, its weighted root-mean-square
variation is 33.1\% of the combined prediction. These ratios describe
the model contribution; the measurement tests the combined potential.

The record spans September 2025 to August 2026, with 54 valid segments
in eight observing periods (Fig. 1c). Its 3,299 ten-minute means contain
545.9 h of accepted nominal one-second measurements. Observations occur
in eight calendar months; the annual span is not continuous operation.
Actual UTC times and all gaps are retained. The June--August
observations supply 273.9 h and sample a broad range of relative
lunar--solar phase, while the earlier periods add different phases and
operating conditions (Fig. 1d). Both the duration and placement of the
observations determine what can be separated.

\hypertarget{jointly-resolved-semidiurnal-responses}{%
\subsubsection{Jointly resolved semidiurnal
responses}\label{jointly-resolved-semidiurnal-responses}}

The reference levels of the 54 observing segments are left independent.
The analysis retains the variation within each segment and uses its
actual UTC times to place that variation at the correct astronomical
phase. All segments can therefore constrain common physical responses
without an artificial joining of their frequency levels. The constants
represent unknown reference levels, not measured frequency jumps. The
full record reveals recurring modulation when ordered by M₂ phase (Fig.
2).

Joint estimation gives \(A_L=1.00\pm0.13\) for M₂+N₂ and
\(A_S=0.93\pm0.36\) for S₂+K₂ (Fig. 3a). Each semidiurnal group has its
own response coefficient within the same physical model. One denotes the
specified combined prediction. The coefficients are estimated from the
original ten-minute means using the same noise weights throughout. The
reported errors are the larger of residual-scaled least-squares and
UTC-day-deletion standard errors; their components and conditional
interval coverage are given in Methods. No empirical daily sinusoid is
subtracted before this joint estimate.

After subtracting the fixed minor tidal remainder, the observations and
unit lunar prediction undergo the same projection against segment
constants and the other two groups. Twelve phase bins reveal the
recurring modulation, with the uncertainty of the jointly estimated
nuisance terms propagated into the error bars. The complete set of
individual observations and the support of each bin are shown below the
phase means.

Allowing a lunar component one quarter-cycle from the predicted phase
gives in-phase and quadrature responses of \(1.00\pm0.13\) and
\(-0.10\pm0.20\) (Fig. 3b), consistent with the astronomical phase.

\hypertarget{spectral-separation-through-the-observation-window}{%
\subsubsection{Spectral separation through the observation
window}\label{spectral-separation-through-the-observation-window}}

The measured spectrum provides a direct view of the two principal
semidiurnal neighbourhoods (Fig. 3c). An equal-weight sinusoidal scan
with 54 segment constants has local maxima at 22.361 and 23.147 µHz,
close to the predicted M₂ and S₂ frequencies of 22.364 and 23.148 µHz.
The prescribed principal frequencies are 0.78 µHz apart. The left axis
expresses spectral density in gravitational-potential units and the
right in fractional-frequency units; a fractional variation of
\(10^{-18}\) corresponds to about \(0.09\,\mathrm{m^2/s^2}\) of
potential difference.

The observation gaps produce sidelobes around a sampled tidal line.
Single-constituent injections through the actual sampling and scan give
8.03\% M₂-to-S₂ cross-power and 8.15\% in the reverse direction.
Conditional scans retain a local feature near each frequency after
projecting out the sine and cosine at the other (Extended Data Fig. 1).
Together with the joint physical-template estimates, these features
identify distinct responses in the neighbouring semidiurnal groups. The
two groups' prescribed internal structure and the sampling-window
response are treated explicitly in Methods.

\hypertarget{consistency-across-observing-periods}{%
\subsubsection{Consistency across observing
periods}\label{consistency-across-observing-periods}}

Observations before and during summer both contribute to the lunar
response. With the same segment model, the coefficient is
\(1.21\pm0.23\) before summer and \(0.77\pm0.18\) in summer. Deleting
each observing period in turn gives point estimates from 0.885 to 1.093
(Fig. 4a). These overlapping fits show each period's influence on the
full-record estimate.

Changing only the grouping of the reference levels gives similar lunar
coefficients (Fig. 4b). The available interval and correction records do
not establish a calibrated transfer of zero across interruptions,
motivating the separate segment constants used here. The eight-period
design gives \(0.987\pm0.123\) on the same observations, physical
templates and weights. Its residual structure and the full comparison of
group coefficients are given in Methods and Extended Data Fig. 2.

Daily temperature changes and operating schedules motivate a 24-h
sensitivity check. Adding one shared 24-h sine--cosine pair to the
present model gives \(A_L=1.04\pm0.13\). The matched summer round-trip
monitor also leaves the lunar estimate essentially unchanged, with a
paired shift of \(0.0061\pm0.0063\) over 246.6 h. The measurement
support and full results of that control are shown in Extended Data Fig.
3. Free 12-h terms are confined to the sensitivity analysis in Methods.

The Shanghai--Wuhan comparison measures tidal variations of the
gravitational potential and separates neighbouring semidiurnal responses
across an interrupted record. The two semidiurnal groups have separately
estimated response coefficients, supported by local peaks near M₂ and S₂
and by the conditional scans. Together they resolve neighbouring lunar
and solar-driven tidal responses within the specified physical grouping.
The phase plot, spectrum and group estimates describe the same
observations from their time, frequency and physical-model perspectives.

Astronomical phase provides the reference that persists between
observing segments. Independent constants accommodate unknown frequency
levels, while the changing lunar--solar phase lets observations at
different dates constrain the neighbouring responses. The result here
concerns separation of neighbouring responses when observations are
interrupted and a common frequency zero has not been transferred between
segments.

This distinction matters for clock networks whose stations cannot
operate continuously. Measurements can retain a common physical phase
reference even when their frequency levels are established separately on
each run. The achievable separation depends on the distribution of
observing dates as well as the accumulated measurement time. Celestial
phase therefore provides a practical reference for resolving adjacent
tidal variations of the gravitational potential in interrupted
optical-clock comparisons. Independent, contemporaneous European
clock-network observations have tracked the overall tidal potential on
different baselines.\textsuperscript{21}

\clearpage

\singlespacing\fontsize{11}{13}\selectfont

\hypertarget{references}{%
\subsection{References}\label{references}}

\begin{enumerate}
\def\labelenumi{\arabic{enumi}.}
\tightlist
\item
  Ludlow, A. D., Boyd, M. M., Ye, J., Peik, E. \& Schmidt, P. O. Optical
  atomic clocks. \emph{Rev.~Mod. Phys.} \textbf{87}, 637--701 (2015).
  \url{https://doi.org/10.1103/revmodphys.87.637}
\item
  Mehlstäubler, T. E., Grosche, G., Lisdat, C., Schmidt, P. O. \&
  Denker, H. Atomic clocks for geodesy. \emph{Rep.~Prog. Phys.}
  \textbf{81}, 064401 (2018).
  \url{https://doi.org/10.1088/1361-6633/aab409}
\item
  Voigt, C., Denker, H. \& Timmen, L. Time-variable gravity potential
  components for optical clock comparisons and the definition of
  international time scales. \emph{Metrologia} \textbf{53}, 1365--1383
  (2016). \url{https://doi.org/10.1088/0026-1394/53/6/1365}
\item
  Denker, H. et al.~Geodetic methods to determine the relativistic
  redshift at the level of \(10^{-18}\) in the context of international
  timescales: a review and practical results. \emph{J. Geod.}
  \textbf{92}, 487--516 (2018).
  \url{https://doi.org/10.1007/s00190-017-1075-1}
\item
  Schröder, S., Stellmer, S. \& Kusche, J. Potential and scientific
  requirements of optical clock networks for validating
  satellite-derived time-variable gravity data. \emph{Geophys. J. Int.}
  \textbf{226}, 764--779 (2021).
  \url{https://academic.oup.com/gji/article/226/2/764/6212226}
\item
  Lisdat, C. et al.~A clock network for geodesy and fundamental science.
  \emph{Nat. Commun.} \textbf{7}, 12443 (2016).
  \url{https://doi.org/10.1038/ncomms12443}
\item
  Riehle, F. Optical clock networks. \emph{Nat. Photon.} \textbf{11},
  25--31 (2017). \url{https://doi.org/10.1038/nphoton.2016.235}
\item
  Chou, C. W., Hume, D. B., Rosenband, T. \& Wineland, D. J. Optical
  Clocks and Relativity. \emph{Science} \textbf{329}, 1630--1633 (2010).
  \url{https://doi.org/10.1126/science.1192720}
\item
  McGrew, W. F. et al.~Atomic clock performance enabling geodesy below
  the centimetre level. \emph{Nature} \textbf{564}, 87--90 (2018).
  \url{https://doi.org/10.1038/s41586-018-0738-2}
\item
  Grotti, J. et al.~Geodesy and metrology with a transportable optical
  clock. \emph{Nat. Phys.} \textbf{14}, 437--441 (2018).
  \url{https://doi.org/10.1038/s41567-017-0042-3}
\item
  Dimarcq, N. et al.~Roadmap towards the redefinition of the second.
  \emph{Metrologia} \textbf{61}, 012001 (2024).
  \url{https://doi.org/10.1088/1681-7575/ad17d2}
\item
  Lindvall, T. et al.~Coordinated international comparisons between
  optical clocks connected via fiber and satellite links. \emph{Optica}
  \textbf{12}, 843--852 (2025).
  \url{https://doi.org/10.1364/OPTICA.561754}
\item
  Hou, L. et al.~Yb/Sr optical clock comparison at \(2.2\times10^{-18}\)
  uncertainty over a 1350-km fiber link. \emph{arXiv}, 2610.34250v1
  (2026). \url{https://arxiv.org/abs/2610.34250v1}
\item
  Jia, Z.-P. et al.~Improved systematic evaluation of a strontium
  optical clock with uncertainty below \(1\times10^{-18}\).
  \emph{Metrologia} \textbf{63}, 025002 (2026).
  \url{https://doi.org/10.1088/1681-7575/ae449e}
\item
  Zhu, Q. et al.~Ytterbium lattice clock with systematic uncertainty of
  \(1.3\times 10^{-18}\) and instability at the \(10^{-19}\) level.
  \emph{arXiv}, 2606.10514v2 (2026).
  \url{https://arxiv.org/abs/2606.10514v2}
\item
  Chen, F.-X. et al.~10⁻²¹-level optical frequency dissemination over
  2067 km of noise-loaded field-deployed fiber network. \emph{Light Sci.
  Appl.} \textbf{15}, 276 (2026).
  \url{https://doi.org/10.1038/s41377-026-02299-1}
\item
  Petit, G. \& Luzum, B. (eds) IERS Conventions (2010). \emph{IERS
  Technical Note} \textbf{36} (2010).
  \url{https://iers-conventions.obspm.fr/content/tn36.pdf}
\item
  CNES. FES2022 (Finite Element Solution) Tidal model (Version 2024).
  CNES \url{https://doi.org/10.24400/527896/A01-2024.004} (2024).
\item
  Farrell, W. E. Deformation of the Earth by surface loads.
  \emph{Rev.~Geophys.} \textbf{10}, 761--797 (1972).
  \url{https://doi.org/10.1029/RG010i003p00761}
\item
  Han, D. \& Wahr, J. The viscoelastic relaxation of a realistically
  stratified earth, and a further analysis of postglacial rebound.
  \emph{Geophys. J. Int.} \textbf{120}, 287--311 (1995).
  \url{https://doi.org/10.1111/j.1365-246x.1995.tb01819.x}
\item
  Stahl, K. et al.~Optical clock network tracks tidal variations of the
  gravity potential. \emph{arXiv}, 2610.39635v1 (2026).
  \url{https://arxiv.org/abs/2610.39635v1}
\end{enumerate}

\clearpage

\singlespacing\fontsize{10}{12}\selectfont

\hypertarget{figure-legends}{%
\subsection{Figure legends}\label{figure-legends}}

\includegraphics{figures/figure1.pdf}

\textbf{Fig. 1 \textbar{} Geometry and changing lunar--solar phase.}
\textbf{a,} Shanghai Sr and Wuhan Yb, 692 km apart, with the 1,350-km
fibre connection through Hefei drawn schematically. \textbf{b,}
Predicted differential M₂ solid and ocean-loading phasors, with the
solid component defining phase zero; amplitudes are fractional clock
rates in units of \(10^{-18}\). \textbf{c,} Actual UTC sampling in eight
observing periods, with accepted nominal member-hours. \textbf{d,} The
S₂−M₂ astronomical phase at those same samples, including nodal
corrections. Summer periods P6--P8 are orange. Geography and phasors
describe the prediction; the lower panels describe the observation
window.

\clearpage

\includegraphics{figures/figure2.pdf}

\textbf{Fig. 2 \textbar{} The full record ordered by lunar phase.}
\textbf{a,} Twelve 30° M₂ phase means of all 3,299 observations, after
subtracting the fixed minor remainder and projecting against 54 segment
constants and the neighbouring semidiurnal and daily groups. Identical
weighted projection and averaging are applied to the unit lunar
prediction; the fitted curve uses the unbinned primary coefficient.
Error bars are one working standard error with the nuisance projection
propagated; phase means are correlated. \textbf{b,} Every projected
observation, with the full observed vertical range. \textbf{c,} Number
of rows in each bin. Data, weights and templates are the primary segment
design throughout. This is a display of the main observations, not an
independent experiment.

\clearpage

\includegraphics{figures/figure3.pdf}

\textbf{Fig. 3 \textbar{} Semidiurnal responses, phase and measured
spectral structure.} \textbf{a,} The two semidiurnal coefficients from
the three-group joint model, using all 3,299 means. Error bars are one
standard error, defined as the larger of scaled least-squares and
UTC-day jackknife errors. The daily group remains in the fit and is
reported in Methods and Extended Data Table 2. \textbf{b,} Lunar
in-phase and quadrature coefficients; the cross marks the unit
prediction. Contours have covariance radii one and two from scaled least
squares, not separately calibrated joint confidence levels. \textbf{c,}
Equal-weight spectral calculation around the principal M₂ and S₂
frequencies, separated by 0.78 µHz: 3,299 real UTC means, 54 segment
constants and equal weights. Blue shows the observed scan; orange
follows subtraction of the fixed unit total prediction. No gap is filled
or compressed. Both axes show the same nominal one-sided projection
density per µHz, related by \(S_W=c^4S_y\). Sampling sidelobes
contribute about 8\% bidirectional M₂/S₂ cross-power. Peak locations
support the principal-frequency attribution; group coefficients in
\textbf{a} use the physical templates and Allan-based weights. All
panels use the same observation rows.

\clearpage

\includegraphics{figures/figure4.pdf}

\textbf{Fig. 4 \textbar{} Consistency across observing periods and
reference-level designs.} \textbf{a,} Lunar estimates from the full
record, the pre-summer and summer subsets, and deletion of each
observing period. Each fit uses segment constants and the primary noise
weights, with errors recalculated on its retained rows. The full,
pre-summer and summer fits contain 3,299, 1,651 and 1,648 means.
Deleting P1--P8 retains 3,136, 3,096, 2,890, 2,852, 2,870, 2,698, 2,704,
2,847 means, respectively. \textbf{b,} Lunar response with 54 segment,
29 operating-group or eight period constants. Only the intercept
grouping changes; all designs use the same 3,299 observations, physical
templates, corrections and noise weights. Error bars are one standard
error using the same larger-of-scaled-least-squares-and-day-jackknife
rule. Shading marks the full-record segment estimate and its standard
error. The estimates share observations; reference-level choice follows
the available calibration records.

\clearpage

\fontsize{12}{14.4}\selectfont\doublespacing

\hypertarget{methods}{%
\subsection{Methods}\label{methods}}

\hypertarget{observations-frequency-sign-and-averaging}{%
\subsubsection{Observations, frequency sign and
averaging}\label{observations-frequency-sign-and-averaging}}

The Shanghai and Wuhan coordinates used by the frozen prediction are
31.130° N, 121.570° E and 30.529° N, 114.357° E. The measured ratio is
Wuhan Yb/Shanghai Sr.~With positive geodetic potential \(W\), its tidal
fractional variation is
\(-[\delta W_{\mathrm{Wuhan}}-\delta W_{\mathrm{Shanghai}}]/c^2\). The
period-relative beat conversion is
\(-d_p(f_{\mathrm{beat}}-f_{{\mathrm{ref}},p})/F_{1550}\), where
\(F_{1550}=1.933992\times10^{14}\) Hz and \(d_p=2\) in 2025 and 1 in
2026. Thus a 2026 beat deviation of 1 Hz has fractional magnitude
\(5.17\times10^{-15}\). The 2025 division factor is applied before
forming the small clock-rate variation.

The supplied valid-interval table defines 54 segments in 29 operating
groups and eight periods. Counter members are read across midnight
before binning. Beijing-time labels are converted to UTC. Summer group
corrections follow the decimal-arithmetic frequency-chain calculation;
early absolute correction entries remain incomplete. Their constant part
is represented by segment intercepts, while any time variation within a
segment remains a possible systematic contribution. Fitted intercepts
are never applied to join the runs for long-term stability estimation.

The membership audit contains 1,966,908 nominal one-second rows before
the minimum-bin-support cut. Three boundary rows reconcile its
difference from the printed interval total. Means are formed in nominal
600-s bins separately within each segment and require at least 180
retained members. The 3,299 accepted means contain 1,965,327 members, or
545.924 h. There are 3,291 unique bin centres: eight centres contain
separate partial means from different segments, with disjoint members.
The 353.39-day span contains observations on 43 UTC dates. The inherited
valid mask is retained; this work neither applies an additional
triangular filter nor selects points using tidal residuals. Exact
raw-file memberships and hashes are retained in the controlled archive.

\hypertarget{what-the-constituent-names-mean}{%
\subsubsection{What the constituent names
mean}\label{what-the-constituent-names-mean}}

Semidiurnal tides produce approximately two cycles per day. M₂ is the
principal lunar constituent; N₂ describes an additional lunar variation
associated with orbital ellipticity. S₂ is the principal solar
semidiurnal constituent, whereas K₂ contains lunar and solar forcing.
The names identify astronomical frequencies, not separate instruments or
independent experiments. The diurnal group contains K₁ (lunisolar), O₁
(principal lunar), P₁ (principal solar) and Q₁ (lunar elliptic). Their
periods are listed in Extended Data Table 1.\textsuperscript{22}
Internal relative amplitudes, phases and nodal modulation are fixed
within each physical group. A group response therefore tests that
complete specified pattern.

\hypertarget{combined-solid-and-loading-potential}{%
\subsubsection{Combined solid and loading
potential}\label{combined-solid-and-loading-potential}}

The physical target includes the direct astronomical potential, the
Earth's induced potential and the effect of clock displacement. The
selected frozen solid sequence implements these through the IERS radial
displacement \(u\) and a fixed effective-potential conversion,

\[y_{\mathrm{solid}}=-\frac{1+k_2-h_2}{h_2}\frac{g(u_{\mathrm{Wuhan}}-u_{\mathrm{Shanghai}})}{c^2},\]

with \(h_2=0.6078\), \(k_2=0.3010\) and \(g=9.794\,\mathrm{m/s^2}\).
This is a displacement-based approximation to the combined solid
potential, not the alternative frequency-resolved direct-potential
implementation. The same conversion also acts on degree-three and
frequency-dependent displacement contributions in the supplied sequence.
A previous implementation comparison found an approximately 18\% K₁
amplitude deficit relative to a potential-weighted reference. The daily
response is accordingly not used to infer a Love number. Uncertainty in
the solid model remains separate from the quoted statistical errors.

For loading, FES2022 sea-surface elevations \(\zeta\) are convolved with
the effective-potential kernel

\[\delta W_{\mathrm{load}}=\int \rho_w\zeta(\mathbf{r}',t)G_W(\psi)\,\mathrm{d}A,\qquad
G_W(\psi)=\frac{ga}{M_E}\sum_n(1+k'_n-h'_n)P_n(\cos\psi).\]

The three terms represent direct water attraction, induced potential and
clock displacement. PREM load Love numbers are retained through degree
696 with the archived high-degree continuation. Eighteen loading
constituents are synthesized with their astronomical nodal corrections.
The complex convention is \(f\,\mathrm{Re}[C\exp(-i\phi)]\); reversing
the exponential sign requires conjugating \(C\). The present analysis
reads the frozen constituent and time-series tables rather than
substituting an older loading sequence.

The eight named constituents form the three physical groups. Minor terms
remain once as a fixed remainder \(r_0\). On the actual weighted
segment-centred sampling, solid, loading and total root-mean-square
variations are respectively 2.928, 1.328 and \(4.011\times10^{-18}\).
Their ratio 0.331 is an amplitude-norm ratio, not an additive fraction
of total power. The separately reported \(0.89\pm0.55\) loading
coefficient belongs to the alternative directly modelled, filtered
reduction. It is unresolved and has not been recalculated as a separate
coefficient of the present estimator.

\hypertarget{joint-estimator-and-zero-point-choice}{%
\subsubsection{Joint estimator and zero-point
choice}\label{joint-estimator-and-zero-point-choice}}

We fit

\[y_i=c_{s(i)}+A_Ls_{L,i}+A_Ss_{S,i}+A_Ds_{D,i}+r_{0,i}+\epsilon_i,\]

with 54 segment intercepts and three physical response coefficients.
Weighted least squares uses the frozen Allan-based uncertainties
\(\sigma_i\), with \(W_{ii}=1/\sigma_i^2\). A period noise coefficient
\(\alpha_p\) is estimated from within-segment overlapping Allan
deviations at 256, 512, 1,024, 2,048 and 4,096 s, after removing the
fixed unit tide. No Allan window crosses a segment boundary. The bin
uncertainty is \(\sigma_i=\alpha_{p(i)}/\sqrt{n_i}\), where \(n_i\) is
its retained member count. This is a working noise model, not stability
calculated after joining fit-adjusted levels. Independent QR replay
reproduces the eight-period design on the same rows before changing only
its intercept grouping. All comparisons therefore hold sampling,
correction values, physical templates and weights fixed.

The daily group is an auxiliary physical component, fitted
simultaneously to propagate its covariance with the semidiurnal
responses. Its coefficient is \(A_D=0.58\pm0.54\) (Extended Data Table
2). Its predicted variation is smaller on this baseline, and its longer
period supplies fewer cycles within short segments, increasing overlap
with segment reference levels and slow variations. This auxiliary
estimate is not a separate measurement objective.

The inspected interval and correction records contain no documented
calibrated zero transfer for the 46 adjacent segment boundaries that
share an eight-period intercept. They give valid times, but not a signed
boundary transfer with uncertainty or a lock/relock history establishing
unchanged reference levels. This motivates separate segment constants.
The same-data residual and deletion diagnostics are reported in Extended
Data Fig. 2 and Table 2. Coefficient proximity to one is not a
model-selection rule.

\hypertarget{statistical-errors-and-conditional-coverage}{%
\subsubsection{Statistical errors and conditional
coverage}\label{statistical-errors-and-conditional-coverage}}

For each group the reported standard error is the larger of
residual-scaled least squares and leave-one-UTC-day-out jackknife
errors. Deletions remove a calendar day, refit all remaining
coefficients and drop empty intercept columns. The jackknife variance is
\((G-1)G^{-1}\sum_g(\hat A_{(-g)}-\overline A)^2\). Period-deletion
errors are recorded separately. Extended Data Table 2 gives all
components and the same-data intercept comparisons. Observations are
repeated measurements of one clock pair, not thousands of independent
experiments.

Conditional simulations apply both designs to the same 3,299-row
sampling and heteroscedastic uncertainties. Each of four scenarios has
1,000 realizations: independent Gaussian noise; equal variance
contributions from Gaussian white noise and a 6-h Ornstein--Uhlenbeck
process within segments; the corresponding 24-h process within periods;
and standardized Student-\(t_5\) noise. The overall scale is the
observed segment-model residual scale; all true physical response
coefficients are one. Each realization recomputes both components of the
standard-error rule. Seed 2026110617 fixes the simulation. Coverage of
\(\hat A\pm1.96u\) is 95.1--97.4\% for the segment-model lunar
coefficient and 95.4--96.8\% for its neighbouring coefficient; Monte
Carlo standard errors are about 0.4--0.7 percentage points (Extended
Data Table 3). These are tests under declared noise laws, not a
calibration against every possible disturbance.

The phase-relaxed model adds the lunar quadrature while retaining the
prescribed M₂/N₂ relationship. Covariance contours use scaled least
squares. Subset and deletion fits refit the same model on their retained
rows. Shared-data estimates are not averaged. The simulation coverage
refers to two-sided intervals; the reported \(\pm\) values are standard
errors.

\hypertarget{phase-means-spectral-density-and-sampling-sidelobes}{%
\subsubsection{Phase means, spectral density and sampling
sidelobes}\label{phase-means-spectral-density-and-sampling-sidelobes}}

For the phase plot, \(N\) contains segment constants and the two
non-lunar groups. We project both \(y-r_0\) and the unit lunar
prediction with \(M_N=I-N(N^TWN)^{-1}N^TW\), and average in twelve fixed
M₂ phase bins. If \(B\) is weighted bin averaging, the mean covariance
is \(BM_NC_yM_N^TB^T\), where \(C_y\) is the scaled working observation
covariance. It retains correlations between means and propagates
nuisance estimation. The full unbinned response remains primary.

Figure 3c preserves the independently specified equal-weight,
54-intercept scan of the unchanged observed means, rather than replacing
it with the amplitude fit's weights. Every trial sine--cosine pair and
the data are projected against the same constants. For a reduction in
squared residuals \(\Delta\mathrm{RSS}(f)\) with dimensionless \(y\),
the nominal one-sided projection density is
\(S_y(f)=600\,\mathrm{s}\,\Delta\mathrm{RSS}(f)\) per Hz, or \(10^{-6}\)
times that value per µHz. Potential density is \(S_W=c^4S_y\); amplitude
density is \(c^2\sqrt{S_y}\). Uniform-grid checks reproduce the
conventional one-sided Fourier density. For interrupted data this
normalization is a stated projection density, not a Parseval-calibrated
variance over a filled year.\textsuperscript{23}

The frequency step is 0.002729 µHz; the inverse calendar span is
0.032752 µHz. Fine plotting steps do not improve physical resolution.
M₂-only and S₂-only astronomical injections use the same actual times,
nodal factors and equal-weight segment projection. Power at the opposite
principal frequency, divided by power at the injected frequency, gives
0.08027 and 0.08149. Conditional scans project out a freely phased
sinusoid at the other frequency before scanning the target
neighbourhood. Amplitude interference means that such corrections act on
sampled waveforms or complex coefficients, not by subtracting scalar
power spectra. Gaps alone do not predict the realization of measurement
noise. Residual density can describe remaining measurement fluctuations,
but includes possible model mismatch and cannot be called a pure
instrumental noise floor.

\hypertarget{instrument-and-daily-response-checks}{%
\subsubsection{Instrument and daily-response
checks}\label{instrument-and-daily-response-checks}}

The matched monitor calculation uses 1,482 rows with 887,882 nominal
members, or 246.634 h, in P6--P8. Both baseline and augmented fits have
identical members, segment labels and weights. Paired day-deletion
replicates propagate uncertainty in their difference. Extended Data Fig.
3 reports the response changes and dates. This comparison concerns the
monitored transmission route during the measured overlap.

On the selected data and model, a common 24-h sine--cosine pair gives
\(A_L=1.036\pm0.133\); separate 24-h pairs in each period give
\(1.003\pm0.136\). These daily sensitivity designs address a periodicity
motivated by temperature and operating schedules, without assigning an
independently measured thermal amplitude. Free 12-h pairs in each period
give \(0.751\pm0.242\) and a very poorly constrained neighbouring
coefficient because they compete with astronomical semidiurnal variation
within short periods. That fit is retained as a sensitivity boundary.

The archived single-bin diagnostic and separately filtered reduction are
documented in Extended Data Table 4. The latter uses different
filtering, weighting, a direct-potential solid model and a
residual-shift uncertainty; its manuscript-reported results are
distinguished from the present replay. Its ERA5 air-temperature
regressors are environmental proxies, distinct from the stated
chamber-thermometer record.

\hypertarget{ai-assisted-preparation}{%
\subsubsection{AI-assisted preparation}\label{ai-assisted-preparation}}

OpenAI Codex assisted English drafting, numerical replay and figure
preparation. {[}Authors to confirm their verification of the analyses,
citations and final text, and responsibility for the submitted work.{]}

\singlespacing\fontsize{11}{13}\selectfont

\hypertarget{data-availability}{%
\subsection{Data availability}\label{data-availability}}

The author review package contains the 3,299 derived means, fixed
templates, noise weights, figure source tables and immutable input
hashes. The original counters and membership ledger are held in the
controlled experimental archive. {[}Authors to supply an
editor-accessible location and confirm the permanent data-access
statement and repository identifier before submission.{]}

\singlespacing\fontsize{11}{13}\selectfont

\hypertarget{code-availability}{%
\subsection{Code availability}\label{code-availability}}

The package contains the independent QR replay, same-data design
comparison, conditional simulations, spectral calculations, covariance
propagation and figure scripts, with software versions and a
reproducible entry point. {[}Authors to confirm the code licence and
permanent repository identifier before submission.{]}

\singlespacing\fontsize{11}{13}\selectfont

\hypertarget{methods-references}{%
\subsection{Methods references}\label{methods-references}}

\begin{enumerate}
\def\labelenumi{\arabic{enumi}.}
\setcounter{enumi}{21}
\tightlist
\item
  NOAA. Harmonic constituents. \emph{NOAA Tides and Currents}.
  \url{https://tidesandcurrents.noaa.gov/harcon.html?id=8577330}
  (accessed 6 November 2026).
\item
  VanderPlas, J. T. Understanding the Lomb--Scargle periodogram.
  \emph{Astrophys. J. Suppl. Ser.} \textbf{236}, 16 (2018).
  \url{https://doi.org/10.3847/1538-4365/aab766}
\end{enumerate}

\clearpage

\singlespacing\fontsize{11}{13}\selectfont

\hypertarget{acknowledgements}{%
\subsection{Acknowledgements}\label{acknowledgements}}

{[}To be completed by the authors.{]}

\singlespacing\fontsize{11}{13}\selectfont

\hypertarget{funding-statement}{%
\subsection{Funding statement}\label{funding-statement}}

{[}Funding sources and grant identifiers to be completed by the
authors.{]}

\singlespacing\fontsize{11}{13}\selectfont

\hypertarget{author-contributions}{%
\subsection{Author contributions}\label{author-contributions}}

{[}To be completed and approved by all authors.{]}

\singlespacing\fontsize{11}{13}\selectfont

\hypertarget{competing-interests}{%
\subsection{Competing interests}\label{competing-interests}}

{[}Declaration to be supplied by the authors.{]}

\singlespacing\fontsize{11}{13}\selectfont

\hypertarget{additional-information}{%
\subsection{Additional information}\label{additional-information}}

\textbf{Correspondence:} {[}Corresponding author names and email
addresses to be supplied; mark them with an asterisk in the author
list.{]}

\clearpage

\singlespacing\sffamily\fontsize{7}{9}\selectfont

\hypertarget{extended-data}{%
\subsection{Extended Data}\label{extended-data}}

\includegraphics{figures/ed_figure2.pdf}

\textbf{Extended Data Fig. 1 \textbar{} Conditional semidiurnal scans.}
Equal-weight scans of the same 3,299 observations with 54 segment
intercepts. \textbf{a,} M₂ neighbourhood before and after projection of
a sine--cosine pair at S₂. \textbf{b,} S₂ neighbourhood before and after
projection at M₂. Dotted lines show the astronomical target frequencies.
These are descriptive projections on shared observations, not additional
detection probabilities.

\clearpage

\includegraphics{figures/ed_figure1.pdf}

\textbf{Extended Data Fig. 2 \textbar{} Same-data comparison of
reference-level designs.} \textbf{a,} Weighted segment-mean residuals
after the eight-period joint fit, ranging from \(-6.75\) to
\(12.00\times10^{-18}\). Segment constants make these means zero by
construction. \textbf{b,} Lunar point estimates after period deletion,
spanning 0.877--1.052 with eight period constants and 0.885--1.093 with
54 segment constants. The weighted residual sums of squares are 3,781.58
and 3,713.64; the latter has 46 additional intercept parameters.
Unweighted residual root-mean-square values are 24.114 and
\(23.934\times10^{-18}\), and reduced residual sums of squares are
1.1501 and 1.1455. Adjacent standardized residual correlations are
−0.127 and −0.148, using pairs within the same segment and separated by
at most 900 s. All comparisons retain identical observations, templates,
corrections and noise weights. These are residual diagnostics, not
measurements of physical jumps or zero-transfer calibration.

\clearpage

\includegraphics{figures/ed_figure3.pdf}

\textbf{Extended Data Fig. 3 \textbar{} Simultaneous summer transmission
control.} Changes caused by adding the fibre round-trip monitor to the
joint model, with paired UTC-day-jackknife standard errors. Baseline and
augmented fits use identical clock/monitor members, labels and weights:
1,482 means containing 887,882 nominal one-second members, or 246.634 h
in P6--P8, from 29 June to 25 August 2026. The lunar, neighbouring and
daily shifts are respectively \(0.00606\pm0.00627\),
\(0.00594\pm0.01453\) and \(0.00791\pm0.01388\). The paired monitor
constrains the measured transmission route on this summer overlap. The
main tidal estimate uses the entire 545.9-h clock record. The prior
transmission-architecture demonstration over 2,067 km is a separate
experiment.\textsuperscript{16}

\clearpage

\textbf{Extended Data Table 1 \textbar{} Astronomical constituent
guide.} Names specify forcing and frequency; one response coefficient
multiplies each complete group with its internal ratios and phases
fixed.

\begin{longtable}[]{@{}lrll@{}}
\toprule\noalign{}
Symbol & Period\_h & Meaning & Group \\
\midrule\noalign{}
\endhead
\bottomrule\noalign{}
\endlastfoot
M2 & 12.421 & Principal lunar semidiurnal & Lunar \\
N2 & 12.658 & Lunar elliptic semidiurnal & Lunar \\
S2 & 12 & Principal solar semidiurnal & Near-12 h \\
K2 & 11.967 & Lunisolar semidiurnal & Near-12 h \\
K1 & 23.934 & Lunisolar diurnal & Diurnal \\
O1 & 25.819 & Principal lunar diurnal & Diurnal \\
P1 & 24.066 & Principal solar diurnal & Diurnal \\
Q1 & 26.868 & Lunar elliptic diurnal & Diurnal \\
\end{longtable}

\clearpage

\textbf{Extended Data Table 2 \textbar{} Estimator and uncertainty
comparison.} All directly replayed designs use the same 3,299 means,
physical templates and Allan-based weights. Errors are standard errors.
The alternative filtered estimate is separately labelled
author-reported.

\begin{longtable}[]{@{}
  >{\raggedright\arraybackslash}p{(\columnwidth - 8\tabcolsep) * \real{0.2706}}
  >{\raggedright\arraybackslash}p{(\columnwidth - 8\tabcolsep) * \real{0.1765}}
  >{\raggedright\arraybackslash}p{(\columnwidth - 8\tabcolsep) * \real{0.1765}}
  >{\raggedright\arraybackslash}p{(\columnwidth - 8\tabcolsep) * \real{0.1765}}
  >{\raggedright\arraybackslash}p{(\columnwidth - 8\tabcolsep) * \real{0.2000}}@{}}
\toprule\noalign{}
\begin{minipage}[b]{\linewidth}\raggedright
Design
\end{minipage} & \begin{minipage}[b]{\linewidth}\raggedright
Lunar
\end{minipage} & \begin{minipage}[b]{\linewidth}\raggedright
Near\_12h
\end{minipage} & \begin{minipage}[b]{\linewidth}\raggedright
Diurnal
\end{minipage} & \begin{minipage}[b]{\linewidth}\raggedright
Status
\end{minipage} \\
\midrule\noalign{}
\endhead
\bottomrule\noalign{}
\endlastfoot
S54 & 1.000 ± 0.132 & 0.925 ± 0.365 & 0.579 ± 0.538 & Replayed \\
G29 & 0.984 ± 0.124 & 1.093 ± 0.298 & 0.759 ± 0.420 & Replayed \\
P8 & 0.987 ± 0.123 & 1.129 ± 0.288 & 0.678 ± 0.359 & Replayed \\
Filtered, 54 segments & 0.97 ± 0.14 & 0.90 ± 0.29 & 0.35 ± 0.51 &
Author-reported \\
\end{longtable}

\begin{longtable}[]{@{}
  >{\raggedright\arraybackslash}p{(\columnwidth - 10\tabcolsep) * \real{0.1486}}
  >{\raggedleft\arraybackslash}p{(\columnwidth - 10\tabcolsep) * \real{0.1622}}
  >{\raggedleft\arraybackslash}p{(\columnwidth - 10\tabcolsep) * \real{0.1757}}
  >{\raggedleft\arraybackslash}p{(\columnwidth - 10\tabcolsep) * \real{0.1351}}
  >{\raggedleft\arraybackslash}p{(\columnwidth - 10\tabcolsep) * \real{0.1757}}
  >{\raggedleft\arraybackslash}p{(\columnwidth - 10\tabcolsep) * \real{0.2027}}@{}}
\toprule\noalign{}
\begin{minipage}[b]{\linewidth}\raggedright
Group
\end{minipage} & \begin{minipage}[b]{\linewidth}\raggedleft
Estimate
\end{minipage} & \begin{minipage}[b]{\linewidth}\raggedleft
Scaled\_LS
\end{minipage} & \begin{minipage}[b]{\linewidth}\raggedleft
Day\_JK
\end{minipage} & \begin{minipage}[b]{\linewidth}\raggedleft
Period\_JK
\end{minipage} & \begin{minipage}[b]{\linewidth}\raggedleft
Reported\_SE
\end{minipage} \\
\midrule\noalign{}
\endhead
\bottomrule\noalign{}
\endlastfoot
Lunar & 1.0004 & 0.132 & 0.1218 & 0.147 & 0.132 \\
Near-12 h & 0.9251 & 0.2993 & 0.3649 & 0.3531 & 0.3649 \\
Diurnal & 0.5792 & 0.3887 & 0.5381 & 0.7735 & 0.5381 \\
\end{longtable}

\clearpage

\textbf{Extended Data Table 3 \textbar{} Conditional 95\% interval
coverage.} Each scenario has 1,000 simulated realizations per design,
with fixed observed sampling and weights. Values are percentages for the
larger-of-LS-and-day-jackknife rule, using multiplier 1.96. They assess
the specified scenarios, not unknown systematic effects.

\begin{longtable}[]{@{}llrrr@{}}
\toprule\noalign{}
scenario & design & Diurnal & Lunar & Near12h \\
\midrule\noalign{}
\endhead
\bottomrule\noalign{}
\endlastfoot
50\% OU 24 h, period & P8 & 97.9 & 98 & 97.3 \\
50\% OU 24 h, period & S54 & 96.4 & 95.1 & 95.4 \\
50\% OU 6 h, segment & P8 & 94.7 & 95.7 & 96 \\
50\% OU 6 h, segment & S54 & 95.8 & 95.6 & 96.3 \\
Student t (df 5) & P8 & 96.8 & 96.3 & 96.2 \\
Student t (df 5) & S54 & 96.2 & 96.1 & 96.8 \\
Gaussian & P8 & 95.7 & 96.2 & 95.9 \\
Gaussian & S54 & 95.5 & 97.4 & 96.7 \\
\end{longtable}

\clearpage

\textbf{Extended Data Table 4 \textbar{} Archived influence diagnostic
and alternative reduction.} These rows identify the observation support
and estimator of each result. The single-bin diagnostic retains its
original eight-period design. The filtered reduction is reported in the
alternative manuscript (v68); its measurements overlap the present data
and its results are not pooled with the primary estimates.

\begin{longtable}[]{@{}
  >{\raggedright\arraybackslash}p{(\columnwidth - 4\tabcolsep) * \real{0.3333}}
  >{\raggedright\arraybackslash}p{(\columnwidth - 4\tabcolsep) * \real{0.3333}}
  >{\raggedright\arraybackslash}p{(\columnwidth - 4\tabcolsep) * \real{0.3333}}@{}}
\toprule\noalign{}
\begin{minipage}[b]{\linewidth}\raggedright
Check
\end{minipage} & \begin{minipage}[b]{\linewidth}\raggedright
Observation support and design
\end{minipage} & \begin{minipage}[b]{\linewidth}\raggedright
Result and status
\end{minipage} \\
\midrule\noalign{}
\endhead
\bottomrule\noalign{}
\endlastfoot
Pre-identified extreme mean & 2026-07-07 01:25 Beijing time, segment 46;
archived eight-period model, fixed templates and weights & Omitting the
mean changes the daily coefficient from 0.678 to 0.721: +0.043 versus
its original 0.359 standard error. The lunar change is +0.039.
Uncertainty was not recalibrated. The mean remains in the present 3,299
rows. \\
Filtered lunar response & Approximately 544 h, 54 segments; 1,200-s
triangular filter every 600 s; equal weights, direct-potential solid
model; continuous-run circular-residual-shift error & \(0.97\pm0.14\),
author-reported. The other grouped results are in Extended Data Table
2. \\
Separate loading coefficient & Same alternative filtered reduction;
loading assigned a separate coefficient & \(0.89\pm0.55\),
author-reported and unresolved. This is not a separate loading estimate
from the primary design. \\
Temperature regressors & Two-site ERA5 2-m air-temperature proxies in
the alternative reduction; distinct from its stated 98-h Sr
chamber-thermometer record & Lunar \(0.98\pm0.13\), author-reported.
Exact regression membership is needed for independent replay. \\
Alternative round-trip control & Reported 474-h overlap in the filtered
reduction & Lunar and whole-template coefficient changes less than
0.005, author-reported. Exact matching memberships are needed for
independent replay; the directly replayed 246.634-h control is shown in
Extended Data Fig. 3. \\
\end{longtable}

The archived influence diagnostic is retained for provenance; the
current observation membership is unchanged.

\end{document}
